Dieses Kapitel beschreibt die Umsetzung der Seite, die Leser:innen tatsächlich zu sehen bekommen: den Feed und das Package, in dem er ausgeliefert wird.\kihinweis{}
Kapitel~\ref{chapter:architecture} hat gezeigt, wo das Package im Gesamtsystem liegt und über welche Schnittstelle es seine Artikel bezieht.
Hier geht es darum, wie es aufgebaut ist, welche Entscheidungen bei der Umsetzung getroffen wurden und welche Probleme dabei aufgetreten sind.
Am Anfang steht der Aufbau des Packages selbst, weil alle weiteren Abschnitte auf seinen Komponenten aufbauen.

\section{Aufbau des Packages}\label{sec:feed:package}
\setauthor{\secondauthor}

Der Diplomarbeitsantrag verlangt, dass der Feed als wiederverwendbares Modul in bestehende Nuxt-Projekte eingebunden werden kann (Anforderungen T4 und E1).
Im Repository liegt der gesamte Feed deshalb nicht in der Feed-App, sondern im Package \texttt{@wb-news/shortform-news} im Ordner \texttt{packages/shortform-news}.
Die Feed-App auf Port 3002 ist selbst nur ein dünner \gls{host}: Sie bindet die Komponente aus dem Package ein und legt fest, was beim Öffnen eines Artikels passiert.
Dadurch wird jede Änderung am Feed automatisch an derselben Stelle wirksam, an der auch fremde Websites sie beziehen.

\subsection{Drei Komponenten mit unterschiedlicher Tiefe}

Websites, die den Feed einbinden, haben unterschiedliche Ausgangslagen.
Eine Website ohne eigene Infrastruktur möchte den fertigen Feed mit möglichst wenig Aufwand übernehmen.
Eine andere Website besitzt vielleicht bereits eine eigene Schnittstelle für Artikel und möchte nur die Darstellung als Karten nutzen.
Das Package bietet daher nicht eine einzige Komponente an, sondern drei, die ineinander verschachtelt sind und sich in ihrem Funktionsumfang unterscheiden (Abbildung~\ref{fig:feed:components}).

\begin{figure}
    \centering
    \begin{tikzpicture}[
        box/.style={draw, rounded corners=2pt, align=left, font=\small, inner sep=6pt},
        lbl/.style={font=\small\bfseries, anchor=north west}
    ]
        \node[box, minimum width=12cm, minimum height=6.4cm, fill=lightgray] (news) at (0,0) {};
        \node[lbl] at ($(news.north west)+(0.15,-0.1)$) {ShortformNewsFeed};
        \node[font=\footnotesize, anchor=north west, align=left, text width=4.2cm] at ($(news.north west)+(7.6,-0.6)$)
            {Laden und Nachladen\\Polling auf Änderungen\\Kategorien, \glqq Für dich\grqq{}\\Likes, Kommentare, Teilen\\Analytics-Ereignisse};

        \node[box, minimum width=7.1cm, minimum height=4.6cm, fill=lightblue] (feed) at (-2.2,-0.55) {};
        \node[lbl] at ($(feed.north west)+(0.15,-0.1)$) {ShortformFeed};
        \node[font=\footnotesize, anchor=north west, align=left, text width=3.0cm] at ($(feed.north west)+(4.0,-0.6)$)
            {Scroll-Container\\Einrasten\\Render-Fenster\\Sichtbarkeit};

        \node[box, minimum width=3.4cm, minimum height=2.8cm, fill=white] (card) at (-3.8,-1.1) {};
        \node[lbl] at ($(card.north west)+(0.15,-0.1)$) {ShortformCard};
        \node[font=\footnotesize, anchor=north west, align=left, text width=3.0cm] at ($(card.north west)+(0.15,-0.6)$)
            {Bild, Titel, \gls{teaser}\\Aktionsleiste\\\glqq Weiterlesen\grqq{}};
    \end{tikzpicture}
    \caption{Verschachtelung der drei Komponenten des Packages und ihre Aufgaben}
    \label{fig:feed:components}
\end{figure}

Die innerste Komponente \texttt{ShortformCard} stellt genau einen Artikel dar.
Sie zeigt das Titelbild, den Titel, den \gls{teaser} und die Aktionsleiste mit Like, Kommentar und Teilen und meldet, wenn jemand den Artikel öffnen möchte.
Eine Karte weiß nichts über andere Karten und lädt selbst keine Daten.

\texttt{ShortformFeed} ordnet beliebig viele Karten untereinander in einem Container an, der beim Wischen auf jeder Karte einrastet.
Diese Komponente entscheidet auch, welche Karten tatsächlich gezeichnet werden, und meldet, wann eine Karte sichtbar wird.
Die Artikel erhält sie als fertige Liste von außen.
Wer eine eigene Datenquelle hat, verwendet diese Komponente und übergibt seine Artikel selbst.

Die äußerste Komponente \texttt{ShortformNewsFeed} ist der vollständige Feed.
Sie benötigt nur die Adresse der Schnittstelle und übernimmt alles Weitere: das seitenweise Laden, die regelmäßige Prüfung auf neue Artikel, die Kategorien, die Sortierung \glqq Für dich\grqq{}, Likes und Kommentare sowie das Senden der Analytics-Ereignisse.
Die Größe der Dateien spiegelt diese Verteilung wider: \texttt{ShortformCard} umfasst 163 Zeilen, \texttt{ShortformFeed} 110 und \texttt{ShortformNewsFeed} 776 Zeilen.
Die Logik ohne Darstellung, etwa das Zusammenführen neuer Daten oder die Berechnung der persönlichen Reihenfolge, liegt in den Dateien \texttt{utils.ts} und \texttt{personalization.ts}.
Sie sind reine TypeScript-Funktionen ohne Bezug zu Vue und lassen sich dadurch ohne Browser mit Unit-Tests prüfen.

Die Aufteilung folgt dem Grundsatz, dass Daten von außen nach innen fließen und Ereignisse von innen nach außen.
Eine Karte meldet ein Like an den Feed, der Feed an \texttt{ShortformNewsFeed}, und erst diese Komponente sendet die Anfrage an die Schnittstelle.
Wie dieser Fluss im Detail umgesetzt ist, behandelt der Abschnitt über Komponenten und Zustand.
Zunächst steht im Vordergrund, wie die Komponenten als Package gebaut und an einen \gls{host} ausgeliefert werden.

\subsection{Build und Auslieferung}

Ein Package muss so gebaut werden, dass ein fremdes Projekt es ohne Kenntnis des Monorepos verwenden kann.
Für den Build kommt unbuild zum Einsatz, ein Build-Werkzeug aus dem UnJS-Ökosystem, auf dem auch Nuxt selbst aufbaut~\cite{unbuild}.
unbuild bietet mehrere Arbeitsweisen an.
Für das Package wird der Modus \texttt{mkdist} verwendet, der jede Quelldatei einzeln in den Ordner \texttt{dist/} überträgt, statt alles zu einer einzigen Datei zusammenzufassen~\cite{mkdist}.
TypeScript-Dateien werden dabei zu JavaScript übersetzt, und zu jeder Datei entsteht eine Typdeklaration.

Vue-Komponenten werden in diesem Modus nicht zu fertigem JavaScript kompiliert.
Die Dateien in \texttt{dist/components} sind weiterhin Single-File Components mit \texttt{<template>}-Block; nur der Skriptteil ist von TypeScript in JavaScript übersetzt.
Das Kompilieren der Templates übernimmt erst der Build des \gls{host}s.
Diese Entscheidung hat zwei Gründe.
Erstens kompiliert der \gls{host} die Komponenten mit seiner eigenen Vue-Version, sodass keine zweite Kopie von Vue im Package mitgeliefert wird.
Zweitens bleiben die Tailwind-Klassen als lesbarer Text in den Templates erhalten, was für die Erzeugung des CSS notwendig ist (siehe unten).
Das fertige \texttt{dist/} ist 116~KB groß, das daraus gepackte Archiv rund 22~KB.

Welche Teile ein \gls{host} importieren kann, legt das Feld \texttt{exports} in der Datei \texttt{package.json} fest.
Tabelle~\ref{tab:feed:exports} zeigt die vier Einstiegspunkte des Packages.

\begin{table}
    \centering
    \small
    \begin{tabularx}{\textwidth}{l >{\raggedright\arraybackslash}X}
        \toprule
        \textbf{Einstiegspunkt} & \textbf{Inhalt} \\
        \midrule
        \texttt{.} & Die drei Komponenten sowie alle Hilfsfunktionen und Typen, etwa der Typ \texttt{Article} \\
        \texttt{./utils} & Nur die Hilfsfunktionen, ohne Vue-Komponenten; genutzt vom Admin-\gls{cms} \\
        \texttt{./nuxt} & Das Nuxt-Modul, das die Komponenten global registriert \\
        \texttt{./tailwind.css} & Eine CSS-Datei, die Tailwind auf die Komponenten des Packages hinweist \\
        \bottomrule
    \end{tabularx}
    \caption{Einstiegspunkte des Packages laut \texttt{package.json}}
    \label{tab:feed:exports}
\end{table}

Vue ist die einzige Laufzeitabhängigkeit des Packages.
Die Bibliothek \texttt{@nuxt/kit}, die nur das Nuxt-Modul benötigt, ist als optionale Peer-Abhängigkeit eingetragen.
Ein Projekt, das Vue ohne Nuxt verwendet, kann die Komponenten dadurch über den Haupteinstiegspunkt (\texttt{.}) importieren, ohne Nuxt installieren zu müssen.

\subsection{Registrierung als Nuxt-Modul}\label{sec:feed:package:module}

In einem Nuxt-Projekt müssen Komponenten aus fremden Packages normalerweise in jeder Datei, in der sie verwendet werden, einzeln importiert werden.
Nuxt erlaubt es Packages jedoch, sich als Modul in den Startvorgang einzuklinken und dabei Komponenten global anzumelden~\cite{nuxtmodules}.
Das Package nutzt diese Möglichkeit mit einem eigenen Modul in der Datei \texttt{src/module.ts}, dessen Kern Listing~\ref{lst:feed:module} zeigt.

\begin{lstlisting}[float=tbp,language=TypeScript,caption={Registrierung der Komponenten im Nuxt-Modul (gekürzt)},label=lst:feed:module]
export default defineNuxtModule<ModuleOptions>({
  meta: {
    name: '@wb-news/shortform-news',
    configKey: 'shortformNews'
  },
  setup() {
    const resolver = createResolver(import.meta.url);
    addComponent({
      name: 'ShortformNewsFeed',
      filePath: resolver.resolve('./components/ShortformNewsFeed.vue')
    });
    // ShortformFeed und ShortformCard analog
  }
});
\end{lstlisting}

Die Funktion \texttt{defineNuxtModule} stammt aus dem Nuxt Kit und beschreibt ein Modul mit Namen, Konfigurationsschlüssel und einer Funktion \texttt{setup}, die beim Start von Nuxt ausgeführt wird~\cite{nuxtmodules}.
Darin meldet \texttt{addComponent} jede der drei Komponenten unter ihrem Namen an~\cite{nuxtkitcomponents}.
Der Pfad wird mit \texttt{createResolver} relativ zur Moduldatei aufgelöst.
Das ist notwendig, weil das Modul nach der Installation im Ordner \texttt{node\_modules} des \gls{host}s liegt und ein fester Pfad dort nicht mehr stimmen würde.

Für einen \gls{host} reduziert sich die Einbindung damit auf einen Eintrag im Feld \texttt{modules} seiner Datei \texttt{nuxt.config.ts}, nämlich den Einstiegspunkt \texttt{@wb-news/\allowbreak{}shortform-news/\allowbreak{}nuxt}.
Danach steht \texttt{<ShortformNewsFeed>} in jeder Seite zur Verfügung, ohne dass ein Import geschrieben werden muss.
Nuxt lädt die Komponenten dabei nur auf Seiten, die sie tatsächlich verwenden.
Der Konfigurationsschlüssel \texttt{shortformNews} ist bereits vorgesehen, wird derzeit aber nicht ausgewertet.
Er ist für spätere Optionen reserviert, etwa ein eigenes Präfix für die Komponentennamen, falls ein \gls{host} bereits Komponenten mit gleichem Namen besitzt.

Die beiden Apps im Monorepo verwenden das Modul nicht.
Sie importieren Komponenten und Hilfsfunktionen direkt, weil sie über den Workspace von pnpm ohnehin Zugriff auf den Quellcode haben.
Ob das Modul in einem fremden Projekt funktioniert, ließ sich im Monorepo deshalb nicht feststellen.
Dafür war ein eigener Test notwendig, den der folgende Abschnitt beschreibt.

\subsection{Einbindung in ein fremdes Nuxt-Projekt}\label{sec:feed:package:host}

Innerhalb des Monorepos funktionierte das Package von Anfang an, weil die Apps auf den Quellcode zugreifen und dieselbe Konfiguration teilen.
Um die Einbindung so zu prüfen, wie sie bei einer echten Website ablaufen würde, wurde im Ordner \texttt{examples/nuxt-host} ein eigenständiges Nuxt-4-Projekt angelegt.
Es liegt bewusst außerhalb des pnpm-Workspace und installiert das Package nicht aus dem Quellcode, sondern als gepacktes Archiv, genau wie ein fremdes Projekt es über eine Registry erhalten würde.
Bei diesem Test traten fünf Fehler auf, die im Monorepo nicht sichtbar gewesen waren.

Der erste Fehler betraf bereits die Installation.
Das Package baute sich ursprünglich im Skript \texttt{postinstall} selbst, das nach jeder Installation ausgeführt wird.
Das Build-Werkzeug unbuild ist aber nur als Entwicklungsabhängigkeit eingetragen und wird bei der Installation in einem fremden Projekt nicht mitinstalliert.
Jede Installation brach dadurch ab.
Die Lösung war, den Build in das Skript \texttt{prepare} zu verschieben.
Laut Dokumentation von npm läuft \texttt{prepare} vor dem Packen eines Packages und bei einer lokalen Installation ohne Argumente, nicht aber, wenn ein Projekt ein Package als Abhängigkeit installiert~\cite{npmscripts}.
Das Archiv enthält damit bereits das fertige \texttt{dist/}, und im Monorepo baut \texttt{pnpm install} das Package weiterhin automatisch.

Der zweite und folgenreichste Fehler betraf das Styling.
Die Komponenten sind ausschließlich mit Klassen von Tailwind CSS gestaltet und bringen kein eigenes, fertiges CSS mit.
Tailwind erzeugt nur für jene Klassen CSS, die es beim Durchsuchen der Projektdateien findet.
Den Ordner \texttt{node\_modules} schließt es dabei grundsätzlich aus~\cite{tailwindsources}.
Im fremden Projekt erschien der Feed deshalb vollständig ohne Gestaltung, obwohl alle Klassen korrekt in den Templates standen.
Abhilfe schafft die Anweisung \texttt{@source}, mit der zusätzliche Pfade für die Suche angemeldet werden~\cite{tailwindsources}.
Damit ein \gls{host} den genauen Pfad in seinen \texttt{node\_modules} nicht kennen muss, liefert das Package die Datei \texttt{tailwind.css} mit (Listing~\ref{lst:feed:tailwind}).
Ihr Pfad ist relativ zur Datei selbst angegeben und stimmt daher unabhängig davon, wo das Package installiert ist.

\begin{lstlisting}[float=tbp,language={},caption={Die mitgelieferte Datei \texttt{tailwind.css} und ihre Verwendung im Host},label=lst:feed:tailwind]
/* packages/shortform-news/src/tailwind.css */
@source './components';

/* CSS-Datei des Hosts */
@import 'tailwindcss';
@import '@wb-news/shortform-news/tailwind.css';
\end{lstlisting}

Die übrigen drei Fehler waren kleiner, zeigen aber dieselbe Ursache: Annahmen, die nur im Monorepo galten.
Eine Animation des Like-Buttons griff auf Einstellungen im Tailwind-Block \texttt{@theme} zurück, die nur die Apps des Monorepos definierten; sie ist jetzt direkt in der Komponente festgelegt.
Die Farbe eines aktiven Likes wurde von einer statischen Klasse für weiße Schrift überdeckt.
In Tailwind v4 entscheidet bei zwei Klassen für dieselbe Eigenschaft die Reihenfolge im erzeugten CSS und nicht die Reihenfolge im \texttt{class}-Attribut; die Farbe wird deshalb jetzt ausschließlich abhängig vom Zustand gesetzt.
Schließlich stellte die Karte jeder Bildadresse einen Schrägstrich voran, wodurch vollständige Adressen auf einen anderen Server ungültig wurden.

Nach diesen Korrekturen zeigt das Beispielprojekt auf Port 3099 denselben Feed wie die Feed-App auf Port 3002, einschließlich Kategorien, Kommentaren und der Übergabe der Artikel-ID an den \gls{host}.
Die Schritte, die ein \gls{host} durchführen muss, sind in der Datei \texttt{README.md} des Packages dokumentiert: das Package installieren, das Modul eintragen, \texttt{tailwind.css} importieren und die eigene Adresse in der Einstellung \texttt{CORS\_ORIGIN} der Schnittstelle freischalten, damit der Browser die Anfragen an die fremde Adresse zulässt (\gls{cors}).
Damit ist die Grundlage beschrieben, auf der alle weiteren Teile des Feeds aufbauen.
Welche Daten zwischen den drei Komponenten fließen und wo der Zustand des Feeds gehalten wird, ist Gegenstand des nächsten Abschnitts.
